\section{Conclusion and Future Work}
\label{sec:conclusion}

%\akshita{slightly rewritten to emphasize contributions}

We develop model ladders and task scaling laws to predict downstream task performance of overtrained LMs (7B--4T and 13B--5T). Specifically, our contributions include predicting \textbf{individual task performance} for a \textbf{range of tasks} with multiple \textbf{design choices}; using a \textbf{fixed set of ladder models} (rather than using a large set of models to find compute-optimal models for each size); and using a small compute budget (\textbf{1\% of target compute}). We also conduct variance analysis to determine task predictability, and provide recommendations for picking good design choices. 

In future work, we hope to explore ways to improve task predictability, such as increasing the size of evaluation sets, alternative evaluation formats, reducing the impact of randomness, etc. We also hope to extend our prediction method to tasks in the multiple-choice (MC) format to more accurately reflect the capabilities of larger LMs. MC accuracy is harder to extrapolate from our smaller ladder models; we show some preliminary results in \S\ref{sec:mc}. Finally, we also hope to validate our method on larger models (70B parameters and beyond).
%\nascomment{beyond (instead of `to') ?} \akshita{extending beyond RC to MC}

% In future work, we hope to reduce the noise in task evaluation metrics by increasing the size of these evaluation sets, which can hopefully yield even lower prediction errors. While we focused on the ranked classification (RC) format of multiple-choice tasks in this work, we hope to make accurate predictions when tasks are written in the multiple-choice (MC) format, which more accurately reflects the capability of larger LMs. MC format accuracy is harder to predict because all our ladder models are too small to exhibit meaningful signal for extrapolation; we show some preliminary results in \S\ref{sec:mc} but leave this as future work. Finally, we also hope to validate our method on larger models (e.g., at the 70B scale and beyond).
